\chapter[Wo sich die Wege zweier Jünglinge]{Wo sich die Wege zweier Jünglinge trennten}
\quellseite{9}

Ein Prediger erzählte einst folgende Geschichte: Vor dreißig Jahren machten sich zwei Jünglinge auf den Weg, um das Park-Theater in der Stadt New York zu besuchen. Es ging ein Spiel über die Bretter, das dazu angethan war, die Religion ins Lächerliche zu ziehen.

Beide Jünglinge waren Söhne frommer Familien. Sie fühlten innerlich, daß sie nicht ins Theater gehen sollten, am allerwenigsten, um dieses Spiel zu sehen, aber dennoch lenkten sie ihre Schritte dahin. So kamen sie schließlich an den Eingang des Schauspielhauses.

Einer der jungen Männer, dem das Gewissen keine Ruhe lassen
\quellseite{10}
wollte, kehrte dort um und machte sich wieder auf den Heimweg. Und obgleich er der Versuchung nicht widerstehen konnte, noch einmal zur Thüre des Theaters zurückzukehren, trug er doch schließlich den Sieg über die Versuchung davon und kehrte der Stätte der Lust den Rücken.

Der andere Jüngling betrat das Theater und folgte einer Versuchung nach der andern. In einem Leben der Sünde und des Leichtsinns sank er, von Stufe zu Stufe. Er verlor seine geschäftliche Stellung, seinen guten Namen, seine unsterbliche Seele und nahm schließlich ein Ende mit Schrecken.

„Ich stehe heute vor euch“, fuhr der Prediger fort, „...und danke Gott, daß ich während der letzten zwanzig Jahre das Evangelium predigen durfte. Ich war der andere Jüngling.“

Es ist nicht auszusagen, wie viele junge Leute jährlich durch den Besuch des Theaters (und ebenso sehr durch den Besuch des Tanzsaals) verderben. Mag es ihnen auch erst als ein unschuldiges Vergnügen vorkommen, so hat es sich doch nur zu oft schon als ein Strick des Satans erwiesen, mit dem er seine Opfer von Christo, von seiner Kirche, vom Glauben und von der ewigen Seligkeit ab- und ins ewige Verderben hineingezogen hat. Ach, daß doch christliche junge Leute sich warnen ließen und von diesen verderblichen Stätten der Weltlust ferne blieben. Gottes Wort sagt: „Wohl dem, der nicht wandelt im Rath der Gottlosen, noch tritt auf den Weg der Sünder, noch sitzt, da die Spötter sitzen; sondern hat Lust zum Gesetz des Herrn und redet von seinem Gesetz Tag und Nacht“, Ps. 1, 1. 2.
